Die einzige Standardableitung auf $\R$  ist $\partial$, deshalb verzichten wir auf einen unteren Index für die Christoffelsymbole. Nach
[[Linearer Zusammenhang/Vertikale Ableitung/Längs Abbildung/Eigenschaften/Fakt|Fakt   (2)]]
ist für

\mavergleichskette
{\vergleichskette
{ \ell,k
}
{ = }{ 1,2
}
{  }{
}
{    }{
}
{   }{
}
}
{}{}{}

\mavergleichskettedisp
{\vergleichskette
{ \tilde{\Gamma}_\ell^k (t)
}
{ =} { \psi_1'(t) \Gamma_{1 \ell}^k (\psi(t)) + \psi'_2(t) \Gamma_{2 \ell}^k (\psi(t))
}
{ =} { \Gamma_{1 \ell}^k (t,e^t)   + e^t \Gamma_{2 \ell}^k (t,e^t)
}
{ } {}
{ } {}
}
{}{}{.}
Mit
[[Halbebene/Poincaré/Riemannsche Geometrie/Christoffelsymbole/Beispiel|Beispiel]]
ergibt sich

\mavergleichskettedisp
{\vergleichskette
{ \tilde{\Gamma}_1^1 (t)
}
{ =} { \Gamma_{1 1}^1 (t,e^t) + e^t \Gamma_{2 1}^1 (t,e^t)
}
{ =} {  -e^t e^{-t}
}
{ =} { -1
}
{ } {
}
}
{}{}{,}

\mavergleichskettedisp
{\vergleichskette
{ \tilde{\Gamma}_1^2 (t)
}
{ =} { \Gamma_{1 1}^2 (t,e^t) + e^t \Gamma_{2 1}^2 (t,e^t)
}
{ =} { e^{-t}
}
{ } {
}
{ } {
}
}
{}{}{,}

\mavergleichskettedisp
{\vergleichskette
{ \tilde{\Gamma}_2^1 (t)
}
{ =} { \Gamma_{1 2}^1 (t,e^t) + e^t \Gamma_{2 2}^1 (t,e^t)
}
{ =} {  -e^{-t}
}
{ } {
}
{ } {
}
}
{}{}{}
und

\mavergleichskettedisp
{\vergleichskette
{ \tilde{\Gamma}_2^2 (t)
}
{ =} { \Gamma_{1 2}^2 (t,e^t) + e^t \Gamma_{2 2}^2 (t,e^t)
}
{ =} { - e^t e^{-t}
}
{ =} { -1
}
{ } {
}
}
{}{}{.}

 [[Kategorie:Latexseite]]
